\documentclass[10pt,a4paper]{amsart}
\usepackage[margin=19mm]{geometry}
\usepackage{amsmath,amssymb,mathtools}
\usepackage[scr=rsfs,cal=euler]{mathalfa}
\usepackage{xfrac}
\usepackage[unicode,hidelinks,bookmarks=false]{hyperref}
\newcommand{\ie}{i.e.\ }
\newcommand{\cf}{cf.\ }
\newcommand{\resp}{respectively\ }
\newcommand{\Subsection}{\subsection}
\providecommand{\nobreakdash}{\nobreak\hbox{-}}
\setlength{\emergencystretch}{2em}
\multlinegap0pt
\usepackage{mdframed}
\usepackage{mdframed}
\usepackage{mdframed}
\usepackage{mdframed}
\usepackage{tikz}
\usepackage{amssymb}
\usepackage{mathtools}
\usepackage{dsfont}
\usepackage{tcolorbox}
\usepackage{enumerate}
\usepackage{float}
\usepackage{mdframed}
\newtheorem{proposition}{Proposition}
\newcommand{\FF}{\mathcal{F}}
\newcommand{\PP}{\mathbb{P}}
\newcommand{\TT}{\mathrm{T}}
\newcommand{\energy}{\mathrm{Energy \,}}
\newcommand{\hh}{\mathfrak{h}}
\title{On the boundedness for the bilinear quadratic functional given by arbitrary strips}
\author{Fr\'ed\'eric Bernicot}
\date{Source: arXiv:2606.22134v1 (CC BY 4.0)}
\begin{document}
\maketitle

\noindent\emph{Source and licence.} On the boundedness for the bilinear quadratic functional given by arbitrary strips. Version arXiv:2606.22134v1 (CC BY 4.0), distributed by arXiv under the Creative Commons Attribution 4.0 International licence, http://creativecommons.org/licenses/by/4.0/, as recorded in the arXiv licence metadata for this exact version and shown on the abstract page https://arxiv.org/abs/2606.22134v1. Full licence notice: \texttt{licenses/arxiv-2606.22134-CC-BY-4.0.txt}.\par\smallskip

\noindent\textit{Editorial scope.} Counted unit: the article's proposition prop:energy2 (source0.tex lines 523--527). The article's complete original proof of it is delivered with it (source0.tex lines 530--568, one \texttt{proof} environment of 39 lines). No mathematics is written, repaired or re-derived: the statement and the proof are the article's own lines, in the article's own notation, and no retained line is substituted or re-flowed.  No further source-local context is needed: every cross-reference of every retained line resolves inside the two delivered spans.  Nothing was omitted from any retained span, so the package carries no omission record.  No retained line contains a citation, so the wrapper carries no bibliography and no bibliography entry.  No retained line contains a figure environment, an image-inclusion command or drawing code, so no image is delivered and the package declares no asset dependency.  Editorial formatting only, in addition: an AMS \texttt{amsart} preamble that restates the article's own declarations and macros used by the retained lines, namely the environments \texttt{proposition} on separate counters, together with the standard packages those lines need.  The retained environments print in this document as Proposition 1 in order of appearance, whereas the article prints the same items with the numbering of its own full document.  The complete native source of the article as distributed in the arXiv e-print for this exact version is bundled unchanged as \texttt{sources/203327/source0.tex}.
\begin{proposition}[{\it Energy} estimates] \label{prop:energy2} Let $\PP'$ be any sub-collection of tri-tiles of $\PP$ and $\hh$ a sequence of functions. Then with $h:=(\sum_{S\in {\mathcal S}} |h_S|^2)^{1/2}$, we have
$$ \energy_{\PP'}(\hh) \lesssim  (\sharp {\mathcal S})^{\alpha} \cdot \|h\|_2.$$
If the collection $\PP'$ is localized on a spatial interval $I_{\PP'}$ then we have the local estimate
$$ \energy_{\PP'}(\hh) \lesssim  (\sharp {\mathcal S})^{\alpha} \cdot \|h \cdot \chi_{I_{\PP'}}\|_2.$$
\end{proposition}

\begin{proof}
Let us chose a maximiser in the definition of the {\it energy}: an integer $n$ and a mutually disjoint forest $\FF$ such that 
$$\energy_{\PP'}(\hh) \leq 2^{n+1}\left(\sum_{\TT \in \FF} |I_{\TT}| \right)^{1/2}$$
and for every tree $\TT\in\FF$
$$ 2^{n r'} |I_{\TT}| \leq \sum_{S\in{\mathcal S}} \sum_{s\in\TT \atop \omega_s \cap S \neq \emptyset} |I_s|^{\alpha r'} |\langle h_S,\phi_{s_3}\rangle|^{r'}.$$
With $c_s:=|I_s|^{\alpha} |\langle h_S, \phi_{s_3}\rangle |$ (where implicitly $S$ is the unique strip in ${\mathcal S}$ intersecting $\omega_{s}$), we have
\begin{align*}
\sum_{\TT\in \FF} \sum_{s\in \TT} |c_s|^{r'} & =
\sum_{\TT\in \FF} \sum_{S\in {\mathcal S}} \sum_{s\in \TT \atop \omega_s \cap S \neq\emptyset} |I_s|^{\alpha r'} |\langle h_S,\phi_{s_3}\rangle|^{r'} \\
& \leq
\sum_{\TT\in \FF} \sum_{S\in {\mathcal S}} \big(\sum_{s\in \TT \atop \omega_s \cap S \neq\emptyset} |I_s|\big)^{\alpha r'} \big(\sum_{s\in \TT \atop \omega_s \cap S \neq\emptyset} |\langle h_S,\phi_{s_3}\rangle|^{2}\big)^{r'/2} \\
& \leq
\sum_{\TT\in \FF} \sum_{S\in {\mathcal S}} |I_{\TT}|^{\alpha r'} \big(\sum_{s\in \TT \atop \omega_s \cap S \neq\emptyset} |\langle h_S,\phi_{s_3}\rangle|^{2}\big)^{r'/2} \\
\end{align*}
where we used H\"older inequality with $\frac{1}{r'}=\alpha+\frac{1}{2}$ and then that for $S\in{\mathcal S}$ beeing fixed, then the frequency tile $\omega_s$ is fixed when $s$ varies along a tree $\TT$ and so the spatial intervals $I_s$ are disjoint (and included in $I_{\TT}$).
So we obtain (by H\"older inequality)
\begin{align*}
\sum_{\TT\in \FF} \sum_{s\in \TT} |c_s|^{r'} & \leq (\sharp {\mathcal S})^{1-\frac{r'}{2}}
\sum_{\TT\in \FF} |I_{\TT}|^{\alpha r'} \big(\sum_{S\in {\mathcal S}} \sum_{s\in \TT \atop \omega_s \cap S \neq\emptyset} |\langle h_S,\phi_{s_3}\rangle|^{2}\big)^{r'/2} \\
& \leq (\sharp {\mathcal S})^{1-\frac{r'}{2}}
\big(\sum_{\TT\in \FF} |I_{\TT}|\big)^{\alpha r'} \big(\sum_{\TT\in \FF} \sum_{S\in {\mathcal S}} \sum_{s\in \TT \atop \omega_s \cap S \neq\emptyset} |\langle h_S,\phi_{s_3}\rangle|^{2}\big)^{r'/2}.
\end{align*}
By easy orthogonality arguments (since for $S\in{\mathcal S}$ fixed, the tri-tiles $s$ such that $\omega_s\cap S\neq \emptyset$ have all the same scale, given by the strip $S$)
\begin{align*}
\big(\sum_{\TT\in \FF} \sum_{S\in {\mathcal S}} \sum_{s\in \TT \atop \omega_s \cap S \neq\emptyset} |\langle h_S,\phi_{s_3}\rangle|^{2}\big)^{1/2} & \lesssim \big(\sum_{S\in {\mathcal S}} \|h_S \|_2^2\big)^{1/2} \\
& \lesssim \|h\|_{2},
\end{align*}
with by convention $h:=(\sum_{S\in {\mathcal S}} |h_S|^2)^{1/2}$. So we conclude to
\begin{align*}
2^{nr'} \big(\sum_{\TT\in \FF} |I_{\TT}| \big) & \leq
\sum_{\TT\in \FF} \sum_{s\in \TT} |c_s|^{r'}  \lesssim (\sharp {\mathcal S})^{1-\frac{r'}{2}}
\big(\sum_{\TT\in \FF} |I_{\TT}|\big)^{\alpha r'} \|h\|_2^{r'} 
\end{align*}
which yields (since $1-\alpha r'=\frac{r'}{2}$)
$$ 2^{n} \big(\sum_{\TT\in \FF} |I_{\TT}| \big)^{\frac{1}{2}} \lesssim (\sharp {\mathcal S})^{\alpha} \cdot \|h\|_2$$
and so allows us to conclude to the (global) energy estimate.

For the localized statement, we just observe that if $\PP'$ is localized on $I_{\PP'}$ then all the previous arguments imply wave-packets also localized in $I_{\PP'}$ and so one can keep and track the localization at every step.
\end{proof}

\end{document}
